\section{Discussion}
\label{sec:discussion}

\paragraph{What the repair establishes.}
Under a design that kept the agent away from the selection and confirmation data, an agent-proposed weight edit
improved this model's handling of false abstention on unseen cases. It met a criterion fixed in advance, and it did
better than three of five simpler interventions under corrected tests. The result is an existence proof under
controlled conditions. It shows that a closed-form, agent-designed rank-1 edit can move this model to a better point
on the recovery--unsupported-answer trade-off than several standard fixes. It does not show that agent-proposed edits
generally do so, or that this edit is the best of its kind.

\paragraph{Why the logit-bias comparison matters.}
A logit bias against ``unknown'' changes nothing inside the model; it only makes one answer less likely at the
output. If an internal edit cannot clearly beat it, the edit's benefit may amount to an output-level shift that
happens to be implemented in the weights. Our data leave this open. The edit gave clearly fewer unsupported answers
than B1 at similar recovery, but the test could not rule out a recovery deficit larger than the 0.05 margin after
correction, and B1 broke no correct answers while the edit broke 3\%. A larger confirmation set, or a comparison
targeted at B1 alone, would be needed to resolve it.

\paragraph{Is there an edit-specific mechanism?}
The evidence does not yet show one. The layer-35 edit was selected and confirmed on behaviour. Its own explanation,
an ``answer available'' signal read by the edit, was not tested in the confirmation; the pathway branch tested the
other submission. Whether the edit's read direction matters, as opposed to its write direction toward cities, is the
first question of the follow-up below.

\paragraph{Improving behaviour is not the same as explaining it.}
The study separated the two on purpose, and they came apart. The agent's two-head edit located a component that is
sufficient to release answers, yet the account the agent gave of that component was not established. The agent's
predictions were often numerically close (31\% of correct answers flipped to ``unknown'' on E, inside its predicted
30--40\%), but close predictions on the agent's own data are weak evidence for an account. A useful edit and a
plausible story can both be produced without the story being right. Evaluations of agent-proposed interventions
should test the two separately.

\paragraph{One model and synthetic tasks.}
Every confirmatory result concerns one 3B model on generated city questions with a fixed instruction and note
format. The false-abstention behaviour itself was much rarer in the three other small models we had examined
(\cref{sec:background}), so the problem this edit solves may be specific to this model and prompt. The T-city check
varies only the answer vocabulary. Natural questions, reworded instructions and misleading contexts
\citep{setiawan2026grabrag} remain untested.

\paragraph{The dose grid.}
The selected dose was the top of its frozen grid, and utility on S was still rising there. The orthogonalisation
comparator was also selected at the top of its grid ($\beta = 1$). A wider grid could change the comparison in either
direction: a larger edit dose might recover more but also break more correct answers. The edit should be read as one
point on a dose curve that was not fully explored.

\paragraph{One agent run.}
The study observed one session of one agent. Another session, another prompt or another model could propose different
edits, better or worse. We did not test whether access to weights was necessary for the agent's success; in a related
auditing setting, a calculation tool mattered more than weight access \citep{liu2026auditors}. Estimating how often
agents produce edits that pass this evaluation would require many runs.

\paragraph{Controls still needed.}
Three controls would sharpen the conclusions: random-head and random-direction controls, to test whether the chosen
components and directions are special; natural-question data with the evidence in view; and comparisons with
component-level methods such as Steer2Edit \citep{steer2edit2026}, weight-folded false-refusal ablation
\citep{wang2025falserefusal}, gated refusal edits \citep{hinkley2026paving} and SAE-feature steering
\citep{hasanov2026drno}. Our comparisons do not establish superiority over editing or steering methods in general.

\subsection{Pre-registered follow-up (in progress; no results reported)}
\label{sec:followup}

A follow-up was pre-registered after the confirmation closed, to ask whether the two edits are specific. It uses a new
controlled split, N (seed 54, 1,950 rows), disjoint from the agent's exploration split and the selection split, and
it cannot change any confirmation verdict above. Its design and manifest were committed before any held-out output
was generated (\cref{app:artefacts}). \textbf{No follow-up results exist at the time of writing, and none are
reported here.}

\begin{itemize}[leftmargin=*]
  \item \emph{Does the layer-35 read direction matter (C4)?} Twenty control edits keep the edit's write vector but use
        a random read vector. Each is scaled on development rows to match the edit's root-mean-square gate and is
        sign-aligned to push in the same direction. C4 passes if the edit's utility exceeds all twenty controls and its
        unsupported-answer rate is below their median. A failure would leave read-side specificity unestablished; on
        its own it would not show equivalence to an output bias.
  \item \emph{Are the two heads unusual (C1, C2)?} C1 requires ablating the target pair to recover more false
        abstentions than each of twenty random pairs from the same layers, by at least 0.30 over their median. C2
        requires doubling the target pair to create more abstentions on previously correct rows than doubling each
        random pair.
  \item \emph{Does the layer-35 signal follow the agent's predictions (C5)?} Four frozen predictions concern how the
        proposed signal changes when evidence is removed, the city changes, distractors are added, or unrelated
        evidence replaces absent evidence.
\end{itemize}
C1, C2 and C4 require an empirical rank of 1 among 21 evaluated edits; C5 uses an intersection--union test at 0.05.
No family-wise error control is claimed. Under a null in which the target's scores are exchangeable with the random
controls' (ties counted against the target), and given a valid C5 test, a union bound puts the chance of at least one
false pass at about 0.19 or less. This bound holds only under that null; it is not a guarantee against every possible
mechanism. Descriptive arms examine single heads, other late-layer pairs, random write vectors, direct logit
attribution and the prompt without its ``unknown'' option. A natural-question check uses 138 LongMemEval development
questions with the evidence in view.
